\chapter{Teorema di Rado}

\section{Teorema}

Sia $<E,F>$ un sistema d'indipendenza vale che $<E,F>$ è un matroide se e solo se $\forall w:E \to R$ Greedy-Max($E,F,w$) risolve il problema di massimo associato a $<E,F>$ e $w$, ossia ritorna un massimale di peso massimo tra tutti i massimali di $F$

\section{Dimostrazione}
Supponiamo che $<E,F>$ è un matroide.

\begin{itemize}
    \item Sia $w:E \to R$ una qualsiasi funzione peso
    \item Sia $S= \{b_1, \dots, b_p\}$ con $w(b_1) \geq \dots \geq w(b_p)$ l'insieme ritornato da Greedy-Max($E,F,w$) sappiamo che \uline{$S \in F$ e $S$ è massimale}
\end{itemize}

Sia ora $A \in F$ con $A$ massimale dobbiamo far vedere che $w(S) \geq w(A)$

Siccome $<E,F>$ è un matroide allora $|A|=|S|=p$ (tutti i massimali hanno la stessa cardinalità)

Quindi possiamo scrivere $A = \{a_1, \dots, a_p\}$ assumendo che $w(a_1) \geq \dots \geq w(a_p)$ mostriamo ora che $\forall i \{1 , \dots, p\}\ w(b_i) \geq w(a_i)$ fatto che consente di concludere che $w(S) \geq w(A)$

Se per assurdo $\exists i \in \{1 , \dots, p\} : w(b_i) < w(a_i)$ Sia $k$ il più piccolo di tali $i$.

Quindi $w(b_k) < w(a_k)$ consideriamo gli insiemi
\begin{itemize}
    \item $D = \{b_1, \dots, b_{k-1}\} \subseteq S \implies D \in F$
    \item $D' = \{a_1, \dots, a_{k-1},a_k\} \subseteq A \implies D' \in F$
\end{itemize}

Siccome $|D'| > |D|$ per la proprietà dello scambio $\exists a_j \in Dì \backslash D _ D \cup \{a_j\} \in F$

Siccome $D$ è costituito dai primi $k-1$ elementi di $S$, quindi i primi $k-1$ elementi selezionati da Greedy-Max e il $k-esimo$ selezionato è $b_k$, essendo $a_j \not \in D, \text{ vale che } w(b_k) \geq w(a_j)$

Inoltre, siccome $a_j \in D'$, vale che $w(a_j) \geq w(a_k)$ quindi da $w(b_k) \geq w(a_j)$ e  $w(b_k) \geq w(a_j)$ si ottiene $w(b_k) \geq w(a_k)$ che contraddice $w(b_k) < w(a_k)$.

Mostriamo che se $<E,F>$ non è un matroide allora $\exists w:E \to R$ tale che Greedy-Max($E,F,w$) non risolve il problema di massimo associato a $<E,F>$ e $w$.

Supponiamo quindi che $<E,F>$ non sia un matroide, allora $\exists A,B \in F$ con $|B|=|A|+1 : \forall b \in B \backslash A\ A \cup \{b\} \not \in F$ sia $w:E \to R$ così definita

{
    \large
    $$
        \forall x \in E, w(x) = 
        \begin{cases}
            1 + \frac{1}{2 \cdot |A|} & \text{se } x \in A \\
            1 & \text{se } x \in B \backslash A \\
            0 & \text{se } x \in (A \cup B)^c
        \end{cases}
    $$
}

Greedy-Max($E,F,w$) ritorna un insieme $S$ nel quale
\begin{itemize}
    \item Vi sono tutti gli elementi di A \\
    (poichè sono quelli di peso > e poichè $A \in F$)

    \item Non vi è alcun elemento di $B \backslash A$\\
    (poichè $\forall b \in B \backslash A\ A \cup \{b\} \not \in F$)

    \item Potrebbero esserci eventualmente elementi di $(A \cup B)^c$\\
    (questi hanno peso nullo)
\end{itemize}

Quindi $w(S)=w(A)=|A| \cdot (1+\frac{1}{2|A|}) = |A|+\frac{1}{2}$. Calcoliamo ora $w(B)$, sapendo che $B=(B \backslash A) \cup (A \cap B)$

\[
    \omega(B) = \underbrace{1 + \dots + 1}_{\text{un } 1 \; \forall x \in B \setminus A} 
    \;+\; 
    \underbrace{1 + \frac{1}{2|A|} + \dots + 1 + \frac{1}{2|A|}}_{\substack{\text{un contributo di} \\ 1 + \frac{1}{2|A|} \; \forall x \in A \cap B}}
\]

Pertanto $w(B) \geq |B| = |A| + 1 > |A| + \frac{1}{2} = w(S)$, in definitiva

{
    \large
    $$
        w(B) > w(S)
    $$
}

Quindi l'insieme S restituito da Greedy-Max non è di peso massimo